\documentclass[10pt,a4paper]{amsart}
\usepackage[margin=19mm]{geometry}
\usepackage{amsmath,amssymb,mathtools}
\usepackage[scr=rsfs,cal=euler]{mathalfa}
\usepackage{xfrac}
\usepackage[unicode,hidelinks,bookmarks=false]{hyperref}
\newcommand{\ie}{i.e.\ }
\newcommand{\cf}{cf.\ }
\newcommand{\resp}{respectively\ }
\newcommand{\Subsection}{\subsection}
\providecommand{\nobreakdash}{\nobreak\hbox{-}}
\setlength{\emergencystretch}{2em}
\multlinegap0pt
\usepackage{bm}
\usepackage{bbm}
\usepackage{dsfont}
\usepackage{xspace}
\usepackage{cancel}
\usepackage{mathtools}
\usepackage{amsthm}
\newtheorem{theorem}{Theorem}
\newtheorem{lemma}[theorem]{Lemma}
\newtheorem*{remark}{Remark}
\newcommand{\OO}{\mathcal{O}}
\renewcommand{\leq}{\leqslant}
\title[Symbols at infinity and Lie algebras of vector fields]{Symbols at infinity and Lie algebras of vector fields}
\author{Emile Bouaziz}
\date{Source: arXiv:2609.25973v1 (CC BY 4.0)}
\begin{document}
\maketitle

\noindent\emph{Source and licence.} Symbols at infinity and Lie algebras of vector fields. Version arXiv:2609.25973v1 (CC BY 4.0), distributed under the Creative Commons Attribution 4.0 International licence, as recorded in the arXiv licence metadata for this exact version and shown on the abstract page https://arxiv.org/abs/2609.25973v1. Full rights notice: licenses/arxiv-2609.25973-CC-BY-4.0.txt.\par\smallskip

\noindent\textit{Editorial scope.} Counted unit: the Lemma whose source label is \texttt{prelim} in the article (source0.tex lines 328--334, source label \texttt{prelim}), delivered together with the article's own complete original proof of it (source0.tex lines 336--350, one \texttt{proof} environment of 15 lines). No mathematics is written, repaired or re-derived: statement and proof are the article's own text line for line. Required context retained uncounted: coord (lines 222--222). Every definition, display and label of those lines is retained unchanged. Omitted from this package and not redistributed: 1--221, 223--327, 335--335, 351--456. Nothing inside a retained excerpt is omitted or altered: every retained body is delivered exactly as the article's lines are written, so no omission receipt of any kind accompanies this package. Citations: the retained excerpts carry no citation command at all, so no bibliography entry of the article is transcribed and no reference is redistributed. Reference adaptation: none was needed and none was made; every reference inside the delivered unit resolves inside it, so no unresolved reference or citation is produced. Printed slips kept literal: no printed word, symbol, hypothesis or number of the counted statement or of its proof was altered. Layout and class: the article's own class is \texttt{amsart} with packages including fontenc, mathpazo, tabularx, hyperref, booktabs, xcolor, tcolorbox, xy, graphicx, tikz-cd, amsthm, amssymb, amsfonts, amsaddr; the wrapper is the lane's pinned \texttt{amsart} editorial wrapper at 10pt on A4, which restates the theorem-like environments the retained text needs and adds the article's own macro definitions that the retained text needs (\texttt{\textbackslash OO}, \texttt{\textbackslash leq}), with extra wrapper packages (bm, bbm, dsfont, xspace, cancel, mathtools). The delivered numbering is the unit's own. Assets: no retained span contains a figure environment, a graphics-inclusion command, a PDF-page inclusion, a table or a TikZ picture; no image file is copied into this package, the package declares no asset dependency and no image file is redistributed with it. Licence: version arXiv:2609.25973v1 of this article is distributed under the Creative Commons Attribution 4.0 International licence, as recorded in the arXiv licence metadata for this exact version and shown on the abstract page https://arxiv.org/abs/2609.25973v1. A full rights notice is delivered at licenses/arxiv-2609.25973-CC-BY-4.0.txt. Characters: every retained excerpt is pure ASCII, so the delivered engine, XeLaTeX with its default Latin Modern text font, reports no missing character.
\begin{remark}\label{coord} Analytic locally we may choose coordinates $(z_1,\cdots,z_n)$ on $X$ so that $D$ and $A$ are given, for some $r\leq n$, by $$D=\{z_1\cdots z_r=0\}\,\,A=\{z_1^{a_1}\cdots z_r^{a_r}=0\}.$$ We will sometimes write $z^B$ for the monomial cutting out the effective divisor $B=\sum_j b_jD_j$, which is to say we set $$z^B=\prod_i z_i^{b_i}.$$ Then $T_X$ is trivialised by the basis $\{\partial_1,\cdots,\partial_n\}$ and $T_X(\log D)$ by $$\{z_1\partial_1,\cdots,z_r\partial_r,\partial_{r+1},\cdots,\partial_n\}.$$ The line bundle $L=\OO_X(A)\subset o_*\OO_{X\setminus D}$ is given by $z^{-A}\OO_X\subset o_*\OO_{X\setminus D}$ and has trivialising section $z^{-A}.$ \end{remark}

\begin{lemma}\label{prelim} The kernel $\mathcal{K}_0$ of the restriction map $\mathcal{S}(B)\to\mathcal{S}(B')$ has a finite length filtration by $\mathcal{S}(B)$-submodules $\mathcal{K}_0\supset\mathcal{K}_1\supset\cdots\supset\mathcal{K}_{a_i}=0$ so that the following two properties are satisfied:
\begin{itemize}
\item The action of $\mathcal{S}(B)$ on the associated graded $\mathrm{Gr}\mathcal{K}$ factors through the restriction morphism $\mathcal{S}(B)\to\mathcal{S}(D_i)$.
\item Considered as modules for the normal subalgebra $\mathcal{N}(D_i)$ and setting $F_c=\iota_i^*\OO_X(-B'-cD_i)$, the associated graded pieces sit in short exact sequences
$$0\to\iota_{i,*}\mathcal{M}_{c,a_i}(F_c)\to\mathrm{Gr}_c\mathcal{K}\to\iota_{i,*}\mathcal{M}_{c,0}(F_c\otimes T_{D_i}(\log\Delta_i))\to0.$$
\end{itemize}
\end{lemma}

\begin{proof} We will construct the filtration globally and then verify its properties locally.


On locally free sheaves on $X$, the kernel of the restriction morphism $\iota_{B,*}\iota_B^*\to\iota_{B',*}\iota_{B'}^*$ is given by tensor product with $$\mathcal{J}_0:=\OO_X(-B')/\OO_X(-B).$$ This has the finite filtration $\mathcal{J}_c:=\OO_X(-B'-cD_i)/\OO_X(-B)$, with associated graded pieces $\mathrm{Gr}_c\mathcal{J}\simeq\iota_{i,*}\iota_i^*\OO_X(-B'-cD_i)=\iota_{i,*}F_c$. Tensoring with $T_X(\log D)\otimes L^{\otimes\bullet}$ gives the required filtration of the kernel, namely $\mathcal{K}_c:=\mathcal{J}_c\otimes T_X(\log D)\otimes L^{\otimes\bullet}$. The projection formula identifies the associated graded pieces as
$$\mathrm{Gr}_c\mathcal{K}\simeq\iota_{i,*}\big(F_c\otimes\iota_i^*T_X(\log D)\otimes(\iota_i^*L)^{\otimes\bullet}\big).$$
The action factors through restriction to $\mathcal{S}(D_i)$ because
$$\big[\OO_X(-D_i)T_X(\log D)\otimes L^{\otimes p},\OO_X(-B'-cD_i)T_X(\log D)\otimes L^{\otimes n}\big]$$
$$\subset\OO_X(-B'-(c+1)D_i)T_X(\log D)\otimes L^{\otimes(p+n)}.$$
Indeed logarithmic vector fields preserve both the ideal sheaf factors.

We turn to the second assertion. On $D_i$ we tensor the normal-tangent sequence with $F_c\otimes(\iota_i^*L)^{\otimes\bullet}$ to give an exact sequence,
$$0\to F_c\otimes(\iota_i^*L)^{\otimes\bullet}e\to F_c\otimes\iota_i^*T_X(\log D)\otimes(\iota_i^*L)^{\otimes\bullet}\to$$ $$\to F_c\otimes T_{D_i}(\log\Delta_i)\otimes(\iota_i^*L)^{\otimes\bullet}\to 0.$$ After pushing forward to $X$ this supplies the relevant exact sequence at the level of the underlying sheaves. 

Now the remaining claims can be proven locally, with respect to coordinates as in remark \ref{coord}. A local section $w$ of $F_c\otimes\iota_i^*L^{\otimes n}e$, of weight $n$, lifts to a section on $X$ of the form $$\widetilde{w}:=z^{B'}z_i^{c}w_0(z) z^{-nA}e,\,\ (w_0(z)\,\,\mathrm{regular}),$$ and a weight $m$ section, $u$, of the normal subalgebra lifts to some $$\widetilde{u}=u_0(z) z^{-mA}e,\,\, (u_0(z)\,\,\mathrm{regular}).$$ Their bracket can be expanded via the Leibniz rule to give $$[\widetilde{u},\widetilde{w}]=[u_0(z)z^{-mA}e,z^{B'}z_i^cw_0(z)z^{-nA}e]$$ $$=(c-a_in+a_im)z^{B'}z_i^cu_0(z)w_0(z)z^{-(n+m)A}e$$ $$+z^{B'}z_i^{c+1}\OO z^{-(n+m)A}e,$$ where the terms of order $z_i^{c+1}$ are the result of both $e(u_0(z))$ and $e(w_0(z))$ being divisible by $z_i$, since $u_0(z)$ and $w_0(z)$ are both regular. The terms of order $z^Bz_i^{c+1}$ dissappear in the associated graded and we deduce that the action is given by $$u\cdot w=(c-a_in+a_im)uw.$$ The kernel term is then a submodule isomorphic to $\mathcal{M}_{c,a_i}(F_c)$, exactly as desired. The cokernel term is dispensed of similarly, let $$\widetilde{w}:=z^{B'}z_i^cw_0(z)vz^{-nA},\,\, (w_0(z)\,\,\mathrm{regular},v\in T_{X}(\log D)),$$ be a lift to $X$ of a weight $n$ local section, $w$, of the extension term,  $$F_c\otimes\iota_i^*T_X(\log D)\otimes\iota_i^*L^{\otimes n}.$$ Then we just compute directly that $$[\widetilde{u},\widetilde{w}]=[u_0(z)z^{-mA}e,z^{B'}z_i^cw_0(z)vz^{-nA}]$$ $$=(c-a_in)z^{B}z_i^cu_0(z)w_0(z)vz^{-(n+m)A}$$ $$+u_0(z)w_0(z)[e,v]z^{-(n+m)A}-v(u_0(z))z^{B'}z_i^cw_0(z)ez^{-(n+m)A}.$$ The second term vanishes upon restriction to $D_i$ as $e$ is central in $T_{D_i}(\log\Delta_i)$, the third term vanishes in the quotient by the normal subspace and in the first term we recognize exactly the module $$\mathcal{M}_{c,0}(F_c\otimes T_{D_i}(\log\Delta_i)).$$
\end{proof}

\end{document}
